\section{Hecke components: an exact cost and a way to avoid it}
\label{sec:hecke}

The source already controls a particular projection onto theta lifts.
Projecting further onto a prescribed Hecke component is a new integral
operation. Its cost is measured by congruences, even when the projector is
orthogonal at the real places. This is the classical setting of congruence
ideals~\cite{hida,ars}; the calculation below makes the cost explicit for
the proposed transfer.

\subsection{A general projector estimate}

\begin{proposition}[Polynomial isolation and denominator]
Let $L$ be a finite free lattice over an integrally closed domain $R$, with
fraction field $F$. Let $U\in\End_R(L)$ be diagonalizable over a field
$K\supset F$, and let its distinct eigenvalues be
$\lambda,\mu_1,\ldots,\mu_s$. Assume the eigenvalues are integral over $R$.
The spectral projector onto the whole $\lambda$-eigenspace is
\[
 e_\lambda=\frac{\prod_{j=1}^s(U-\mu_j I)}
                    {\prod_{j=1}^s(\lambda-\mu_j)}.
\]
In an integral coefficient ring containing the eigenvalues, the denominator
$\delta_\lambda=\prod_j(\lambda-\mu_j)$ clears this projector. At a
nonarchimedean valuation $v$ nonnegative on that integral coefficient ring,
its entries have valuation at least
$-v(\delta_\lambda)$.
\end{proposition}
\begin{proof}
The numerator acts on the $\lambda$-eigenspace by $\delta_\lambda$, and on
each other eigenspace by zero. It is an integral endomorphism. Division
therefore gives the projector and the valuation bound. Repeated
eigenvalues need not be repeated in the product, because $U$ is semisimple.
\end{proof}

This is an upper bound, not a claim that every such denominator is minimal.
A block projector can often be obtained with a smaller denominator, or no
denominator. Conversely, a congruence between the chosen eigencharacter and
its complement can make an integral splitting impossible. An archimedean
orthogonal projector has norm one but can have nonintegral entries in the
arithmetic lattice. These two facts are compatible.

\subsection{The complete weight-24 calculation}

Let $q=e^{2\pi i z}$, and put
\[
 E_4=1+240\sum_{n\ge1}\sigma_3(n)q^n,
 \qquad \Delta=q\prod_{n\ge1}(1-q^n)^{24}.
\]
The usual dimension formula gives
$\dim_\C S_{24}(\SL_2(\Z))=2$~\cite{hida}.
The integral cusp lattice has basis
\[
 g_1=\Delta E_4^3-696\Delta^2=q+0q^2+\cdots,
 \qquad g_2=\Delta^2=q^2+\cdots.
\]
Indeed these independent integral forms span over $\C$, and the first two
coefficients of any integral form are exactly its coordinates in this
basis. The formula
\[
 (T_p f)_n=a_{pn}+p^{23}a_{n/p},
 \qquad a_{n/p}=0\quad(p\nmid n),
\]
then gives
\begin{equation}\label{eq:t2}
 U=T_2=\begin{pmatrix}0&1\\20468736&1080\end{pmatrix},
 \qquad \det(XI-U)=X^2-1080X-20468736.
\end{equation}
The discriminant is
\[
 83041344=576\cdot144169=2^6 3^2\cdot144169.
\]
The integer $144169$ is prime and is $1$ modulo $4$. Set
$s=\sqrt{144169}$, $K=\Q(s)$, and
$\OO_K=\Z[(1+s)/2]$. The eigenvalues are
$\lambda_\pm=540\pm12s$, and the two normalized eigenforms are
$g_1+\lambda_\pm g_2$.

\begin{proposition}[Exact projector denominator in weight 24]
On $L\otimes\OO_K$, the denominator ideal of either eigenprojector is
\[
 \{a\in\OO_K:a e_\pm(L\otimes\OO_K)\subseteq L\otimes\OO_K\}
       =(24s).
\]
Its norm is $576\cdot144169=83041344$. The smallest positive rational
integer $d$ satisfying $d e_\pm(L\otimes\OO_K)\subseteq L\otimes\OO_K$
is
\[
 d=24\cdot144169=3460056.
\]
\end{proposition}
\begin{proof}
We have
\[
 e_+=\frac1{24s}
 \begin{pmatrix}-540+12s&1\\20468736&540+12s\end{pmatrix},
 \qquad e_-=I-e_+.
\]
The numerator is integral, so $(24s)$ clears the entries. Its upper-right
entry is exactly $1/(24s)$; hence every clearing scalar must belong to
$(24s)$. For a rational integer $d$, the condition
$d/(24s)=(d/(24\cdot144169))s\in\OO_K$ implies
$d/(24\cdot144169)\in\Z$: writing $cs=2c(1+s)/2-c$ shows that both
coordinates are integral exactly when $c\in\Z$. This proves minimality.
The displayed value of $d$ clears every entry, and the ideal norm follows
by multiplying the two conjugates of $24s$ and taking absolute value.
\end{proof}

The result cannot be improved by choosing another combination of Hecke
operators while demanding the same eigenspace and complement: the
projector is unique. In fact, every integral matrix commuting with $U$ has
the form
\[
 aI+bU=\begin{pmatrix}a&b\\20468736b&a+1080b\end{pmatrix}
 \quad(a,b\in\Z).
\]
This follows by multiplying two $2\times2$ matrices and equating entries.
Every $T_n$ is integral and commutes with $T_2$, so the full integral Hecke
algebra here is exactly $\Z[I,U]$; the reverse inclusion uses $T_1=I$ and
$T_2=U$. This also explains the congruence primes $2,3,144169$ seen in the
Hecke algebra discriminant. The example and congruence phenomenon are
classical; the present contribution is their explicit use in costing a
proposed transfer, with a reproducible certificate.
Writing $\omega=(1+s)/2$ gives $\lambda_+=528+24\omega$; the Hecke
order is $\Z+24\OO_K$, of conductor $24$ in $\OO_K$.

\subsection{Keep the whole Galois orbit when the goal permits it}

\begin{proposition}[Entire products avoid eigenform isolation]
Let $L_1,\ldots,L_t$ be holomorphic functions, none identically zero, on the critical
strip $0<\Re s<1$, and let $L_{\mathrm{all}}=\prod_j L_j$.
Every zero of $L_{\mathrm{all}}$ lies on $\Re s=1/2$ if and only if that
is true of every $L_j$. At each point $\rho$,
\[
 \operatorname{ord}_{\rho}L_{\mathrm{all}}
      =\sum_j\operatorname{ord}_{\rho}L_j,
\]
and all summands are nonnegative.
\end{proposition}
\begin{proof}
Factor each holomorphic function locally as
$(s-\rho)^{m_j}u_j(s)$, with $u_j(\rho)\ne0$ and $m_j\ge0$.
Multiplying gives the asserted order and both implications.
\end{proof}

For a Galois orbit of holomorphic eigenforms of the same weight and trivial
nebentypus, normalize all $L$-functions to center $1/2$ and multiply them.
They are entire in the strip, and their prime traces add. If the orbit has
$t$ members, Deligne's bounds give a nonnegative deficit
\[
 2t-\frac{\tr(T_p\mid\text{orbit})}{p^{(k-1)/2}}\ge0.
\]
This is the orbit version of the analytic trace-deficit target already
present in the earlier report~\cite{deficit,deligne}. In weight 24 the full
space is one orbit, so its projector is simply $I$. No conjecture about
Hecke irreducibility in general weights is needed for this example. One
could instead take an entire fixed-weight space, provided a prospective
arithmetic argument controls that whole space.
If a larger space has several Galois orbits, projecting onto one entire
orbit can still incur denominators from congruences with the other orbits.
Only the cost of separating the conjugates within that orbit is avoided;
in the weight-24 example that orbit is the whole space.

This is a useful simplification of the target, not a proof of its RH.
The source still does not realize these trace coefficients in its theta
divisibility counts, and a whole-orbit formulation does not itself solve
that problem. The half-integral unitary normalization also requires an
honest comparison of algebraic and archimedean factors; it cannot be used
as an unexplained nonintegral divisibility exponent. For meromorphic
Artin factors, zeros can cancel poles, so the proposition cannot be used
without its holomorphy hypothesis.
